Una masa de \(m=100\,\mathrm{kg}\) cuelga de un resorte de rigidez \(k=80\,\mathrm{N/m}\) y un amortiguador viscoso \(c=40\,\mathrm{N\,s/m}\), en paralelo y fijados a un soporte inmóvil. Se mide x hacia abajo desde el equilibrio previo bajo la gravedad. Desde \(t=0\) actúa \(F(t)=F_0\sin(2\pi ft)\), con \(F_0>0\) en newtons, \(f=2\,\mathrm{Hz}\), y \(x(0)=x'(0)=0\). Se supone respuesta lineal y parámetros constantes.

\begin{enumerate}[label=(\alph*),ref=\alph*,leftmargin=*]
\item Formula la ecuación diferencial.
\item Obtén x(t).
\item Dibuja la respuesta.
\item ¿Qué sucedería si una frecuencia real f satisficiera
\[2\pi f=\frac{\sqrt{c^2-4km}}{2m}\ ?\]
\end{enumerate}
